\documentclass[a4paper,11pt]{article}
\usepackage[utf8]{inputenc}
\usepackage[T1]{fontenc}
\usepackage[french]{babel}
\usepackage{lmodern}
\usepackage{amsmath,amssymb,amsthm}
\usepackage{mathrsfs}
\usepackage{enumitem}
\usepackage{multicol}

\setlist[enumerate]{label=\textbf{\textcolor{Couleur8}{\arabic*.}}, left=1em}

% Si vous voulez aussi modifier les listes enumerate imbriquées :
\setlist[enumerate,2]{label=\textcolor{Couleur8}{\alph*)}}
\setlist[enumerate,3]{label=\textcolor{Couleur8}{\roman*)}}
\setlist[enumerate,4]{label=\textcolor{Couleur8}{\Alph*)}}
\usepackage{graphicx}
\usepackage{xcolor}
\usepackage{tcolorbox}
\usepackage{geometry}
\usepackage{fancyhdr}
\usepackage{titlesec}
\usepackage{cancel}
\usepackage{ifthen}
\usepackage{tikz}
\usetikzlibrary{arrows.meta}
\usepackage{tkz-tab}
\usepackage{eurosym}

% OPTION PROF/ÉLÈVE
% Décommentez UNE des deux lignes suivantes :
\newboolean{prof}\setboolean{prof}{true}   % Version professeur (avec corrections des devoirs)
\newboolean{prof}\setboolean{prof}{true}  % Version élève (sans corrections des devoirs)

\definecolor{Couleur1}{HTML}{4A5D7A} %Th
\definecolor{Couleur2}{HTML}{A5B5D6} %Prop
\definecolor{Couleur3}{HTML}{A8C68A} %Méthode
\definecolor{Couleur6}{HTML}{8A9CC1} %Def
\definecolor{Couleur7}{HTML}{E6B459} %Rq Voc
\definecolor{Couleur8}{HTML}{D36740} %Titres
\definecolor{correction}{HTML}{D36740} % Couleur pour les corrections
\definecolor{coopmathscorr}{HTML}{F15929} % Couleur corrections coopmaths

% Configuration des marges
\geometry{
	top=2cm,
	bottom=2cm,
	inner=1.5cm,
	outer=1.5cm,
}

% Police
\renewcommand{\familydefault}{\sfdefault}

% Compteurs
\newcounter{seancenum}
\newcounter{seancenumD}
\setcounter{seancenum}{1}
\setcounter{seancenumD}{1}

% Configuration des en-têtes et pieds de page
\pagestyle{fancy}
\fancyhf{}
\ifthenelse{\boolean{prof}}{
    \fancyhead[L]{\textcolor{Couleur8}{\textbf{Corrections Automatismes - Classe de seconde - Version Professeur}}}
}{
    \fancyhead[L]{\textcolor{Couleur8}{\textbf{Corrections Devoirs - Séance 14 - Classe de seconde}}}
}
\fancyhead[R]{\textcolor{Couleur8}{\textbf{2025-2026}}}
\fancyfoot[L]{\textcolor{Couleur8}{Mme Bertail et Mme Pinto}}
\fancyfoot[R]{\textcolor{Couleur8}{\thepage}}
\renewcommand{\headrulewidth}{1pt}
\renewcommand{\headrule}{\hbox to\headwidth{\color{Couleur8}\leaders\hrule height \headrulewidth\hfill}}

% Environnements personnalisés
\newtcolorbox{seance}[1]{
    colback=Couleur6!10,
    colframe=Couleur1!65,
    fonttitle=\bfseries\large,
    title=Correction Séance \theseancenum #1,
    left=5mm,
    right=5mm,
    top=3mm,
    bottom=3mm,
    arc=0mm,
    after=\stepcounter{seancenum}
}

\newtcolorbox{seanceD}[1]{
    colback=Couleur6!5,
    colframe=Couleur6!45,
    fonttitle=\bfseries\large,
    title=Correction Devoirs - Séance \theseancenumD #1,
    left=5mm,
    right=5mm,
    top=3mm,
    bottom=3mm,
    arc=0mm,
    after=\stepcounter{seancenumD}
}

\begin{document}

\setcounter{seancenumD}{14}
\newpage
\begin{seanceD}{} %14

\begin{enumerate}
\item On reconnaît une équation produit-nul, donc on applique la propriété :\\
{\color{correction}Un produit est nul si et seulement si au moins un de ses facteurs est nul.}\\
$(x-4)(7x+5)=0$\\
$x-4=0$ ou $7x+5=0$\\
$ x=4$ ou $7x=-5$\\
$ x=4$ ou $x=\dfrac{-5}{7}$\\
On en déduit : $S={\color{correction}\boldsymbol{\left\{-\dfrac{5}{7};4\right\}}}$.

\item On met l'expression au même dénominateur :\\
$\begin{aligned}
\dfrac{1}{2}-\dfrac{3x+3}{x}&=\dfrac{x-2\times \left(3x+3\right)}{2x}\\
&=\dfrac{x -6x -6}{2x}\\
&=\dfrac{-5x -6}{2x}\\
\end{aligned}$\\
$\phantom{\dfrac{1}{2}-\dfrac{3x+3}{x}}=-\dfrac{5x +6}{2x}$\\


\item Déterminer les antécédents de $2$ revient à déterminer les nombres qui ont pour image $2$.\\
On part de $2$ sur l'axe des ordonnées et on lit les antécédents (éventuels) sur l'axe des abscisses.\\
On en trouve $2$ : ${\color{correction}\boldsymbol{3\,;\,5}}$.

\item On reconnaît le développement de l'égalité remarquable :\\
$(a-b)^2=a^2-2ab+b^2$ avec $a=x$ et $b=2$.\\
On a donc :\\
$-4x+4+x^2={\color{correction}\boldsymbol{(x-2)^2}}$

\item $\bullet$ La seule égalité correcte est : ${\color{correction}\boldsymbol{\dfrac{\dfrac{3}{4}}{5}= \dfrac{3}{20}}}$.\\
$\dfrac{\dfrac{3}{4}}{5}=\dfrac{3}{4}\times \dfrac{1}{5}=\dfrac{3}{4\times 5}=\dfrac{3}{20}$\\
$\bullet$ L'égalité $\dfrac{\dfrac{5}{6}}{5}=\dfrac{25}{6}$ est fausse.\\
$\dfrac{\dfrac{5}{6}}{5}=\dfrac{5}{6\times 5}=\dfrac{5}{30}$\\
$\bullet$ L'égalité $\dfrac{5}{6}=\dfrac{25}{36}$ est fausse.\\
$\dfrac{25}{36}=\left(\dfrac{5}{6}\right)^2\neq \dfrac{5}{6}$\\
$\bullet$ L'égalité $5\div \dfrac{5}{6}=\dfrac{5}{30}$ est fausse.\\
$5\div \dfrac{5}{6}=5\times \dfrac{6}{5}=\dfrac{5\times 6}{5}=\dfrac{30}{5}$\\

La bonne réponse est la réponse {\bfseries \color{correction}B}.

\end{enumerate}

\end{seanceD}

\end{document}